% !TEX root = paper.tex
\newcommand{\papername}{TruthTable++}
\newcommand{\truthtable}{\mathsf{TruthTable}}

% FORMATTING
\allowdisplaybreaks
\newcommand{\FormatAuthor}[3]{
\begin{tabular}{c}
#1 \\ {\small\texttt{#2}} \\ {\small #3}
\end{tabular}
}
\newcommand{\doclearpage}{%
\iflncs
\clearpage
\fi
}
\newcommand{\DoQuote}[1]{``#1''}
\newcommand{\defemph}[1]{\textbf{\emph{#1}}}
\newcommand{\keywords}[1]{\bigskip\par\noindent{\footnotesize\textbf{Keywords\/}: #1}}
\providecommand{\CryptoEprint}[1]{Cryptology ePrint Archive, Report}



\newcounter{piopnum}
\setcounter{piopnum}{0}
\crefname{PIOP}{PIOP}{PIOPs}
\crefformat{PIOP}{#2#1#3}
\Crefformat{PIOP}{#2#1#3}

\makeatletter
% Args: * (optional star) forces the box to stay on one page/column;
%       #2 full title (shown in the box), #3 label key (piop:#3),
%       #4 optional short name used as the \cref link text (defaults to the title).
\NewDocumentEnvironment{DefinePIOP}{s m m O{#2}}{%
  % Default nobreak=false lets a long box split across columns (two-column S&P)
  % and across pages; short boxes still stay whole since a break only happens
  % when the box does not fit. The star sets nobreak=true to keep the box
  % together (it will overflow the page bottom if taller than the text height).
  \IfBooleanTF{#1}{\def\dpiop@nobreak{true}}{\def\dpiop@nobreak{false}}%
  \begin{mdframed}[nobreak=\dpiop@nobreak, splittopskip=\topskip, splitbottomskip=0.4em]
    \small
    \begin{center}
      \refstepcounter{piopnum}%
      \textbf{PIOP \thepiopnum: #2}%
      \protected@edef\@currentlabelname{#4}%
      \def\@currentlabel{#4}%
      \protected@edef\cref@currentlabel{%
        [piopnum][\arabic{piopnum}][]\unexpanded{#4}}%
      \let\@currentcounter\@empty
      \label[PIOP]{piop:#3}%
    \end{center}
    \vspace{0.5em}
    \noindent
}{%
  \end{mdframed}
}
\makeatother

% ---- Algorithm boxes (adapted from the int-ext project) ----
\newcounter{algnum}
\setcounter{algnum}{0}
\crefname{Alg}{Algorithm}{Algorithms}
\Crefname{Alg}{Algorithm}{Algorithms}
\makeatletter
% Args: #1 title (shown in the box), #2 label key (alg:#2), #3 body.
\newcommand{\DefineAlg}[3]{%
  \begin{mdframed}[nobreak=true,leftline=false,rightline=false]
    \small
    \begin{center}
      \vspace{0.5em}
      \refstepcounter{algnum}%
      \textbf{Algorithm \thealgnum:  #1}
      \rule{\dimexpr\linewidth-\fboxsep-\fboxrule\relax}{0.4pt}
      \protected@edef\@currentlabelname{#1}%
      \label[Alg]{alg:#2}%
    \end{center}
    \noindent
    #3
  \end{mdframed}
}
\makeatother


%%%%%%%%
\newcommand{\F}{\mathbb{F}}
\newcommand{\eqPoly}{\mathsf{Eq}}
\newcommand{\contig}{\mathcal{C}}
\newcommand{\contigmus}{\contig_{\mu, s}}

%%%%%%% 
\newcommand{\Prover}{\mathbf{P}}
\newcommand{\Verifier}{\mathbf{V}}
\newcommand{\Vector}[1]{\mathbf{#1}}



%%%%%%%
\DeclarePairedDelimiter\DoubleBrackets{\llbracket}{\rrbracket}
\newcommand{\Oracle}[1]{\DoubleBrackets{#1}}


% NP relation = (index, instance, witness)
\newcommand{\NPRelation}{\mathscr{R}}
\newcommand{\NPLanguage}{\mathscr{L}}
\newcommand{\NPParams}{\mathbbm{p}}
\newcommand{\NPIndex}{\mathbbm{i}}
\newcommand{\NPInstance}{\mathbbm{x}}
\newcommand{\NPInstancey}{\mathbbm{y}}
\newcommand{\NPInstances}{\bm{x}}
\newcommand{\NPInstanceSizeBound}{\mathsf{N}}
\newcommand{\NPRelationWithSizeBound}{\NPRelation_{\NPInstanceSizeBound}}
\newcommand{\NPWitness}{\mathbbm{w}}
\newcommand{\ValidNPInstance}{\NPLanguage(\NPRelation)}
\newcommand{\ValidNPInstanceWithSizeBound}{\NPLanguage(\NPRelationWithSizeBound)}
\newcommand{\Yes}[1]{#1^{{\scriptscriptstyle\mathrm{YES}}}}
\newcommand{\No}[1]{#1^{{\scriptscriptstyle\mathrm{NO}}}}

\newcommand{\actset}{A}
\newcommand{\act}{\mathsf{a}}
\newcommand{\data}{\mathsf{d}}
\newcommand{\col}{\mathsf{c}}
\newcommand{\MLE}[1]{\tilde{#1}}

\newcommand{\Relation}[1]{\mathcal{R}_{#1}}
\newcommand{\Language}[1]{\mathcal{L}_{#1}}
\newcommand{\op}{\mathsf{op}}

\newcommand{\PIOPProver}{\bf{P}}
\newcommand{\PIOPVerifier}{\bf{V}}
\newcommand{\idx}{\mathcal{I}}
\newcommand{\idxrot}{\idx^{\mathsf{rot}}_\mu}
\newcommand{\Table}{\mathsf{T}}
\newcommand{\BoolHCube}[1]{\mathscr{H}_{#1}}
\newcommand{\TableIn}{{\mathsf{I}}}
\newcommand{\TableOut}{{\mathsf{O}}}
\newcommand{\TableLeft}{{\mathsf{L}}}
\newcommand{\TableRight}{{\mathsf{R}}}
% Low-level PIOP / keyed-sumcheck helpers
\newcommand{\samples}{\overset{\$}{\leftarrow}}
\newcommand{\mult}{\cdot}
\newcommand{\key}{\mathsf{key}}
\newcommand{\val}{\mathsf{val}}
\newcommand{\Left}{\mathsf{L}}
\newcommand{\Right}{\mathsf{R}}
\newcommand{\elem}{\kappa}
\newcommand{\lookup}{\sqsubseteq}

% Author notes: shown only when \ifnotes (see configs/preamble.tex).
% Usage: \ali{this needs a proof}, \pratyush{rephrase}, etc.
\newcommand{\authnote}[3]{\ifnotes{\normalfont\sffamily\bfseries\color{#2}[#1: #3]}\fi}
\newcommand{\ali}[1]{\authnote{Ali}{cyan}{#1}}
\newcommand{\bharath}[1]{\authnote{Bharath}{blue}{#1}}
\newcommand{\pratyush}[1]{\authnote{Pratyush}{green!50!black}{#1}}
\newcommand{\ryan}[1]{\authnote{Ryan}{orange!85!black}{#1}}






% Polynomial Commitment Schemes
\newcommand{\Commitment}{\mathsf{cm}}